% ======================================================== evaluation ======
\section{Reference implementation and evaluation}
\label{sec:eval}

\subsection{What was built}

Roughly 3{,}000 lines of dependency-light Python implementing the decision core
and the infrastructure patterns described above, plus a discrete-event
simulation harness. The harness drives the \emph{real} engine \textemdash{}
\texttt{core.optimizer}, \texttt{core.priority}, \texttt{core.matching},
\texttt{infra.store} \textemdash{} so the numbers below describe the actual
system rather than a model of one.

\begin{center}\footnotesize
\begin{tabular}{@{}lp{80mm}@{}}
\toprule
\textbf{Module} & \textbf{Contents} \\
\midrule
\texttt{domain/models} & Entities, state machines, assignment maturity, journey projection \\
\texttt{core/geo} & Haversine, cell indexing, bucketed spatial index with expanding-ring search \\
\texttt{core/eta} & Travel-time service: versioned cache, circuit breaker, pessimistic great-circle fallback, road-condition overlay \\
\texttt{core/priority} & Explainable additive triage model with bounded aging \\
\texttt{core/matching} & Jonker--Volgenant exact assignment, greedy warm start, LNS \\
\texttt{core/coverage} & Self-exciting demand forecast, reserve protection, idle relocation \\
\texttt{core/optimizer} & Two-tier engine, geographic sharding, churn control, decision records \\
\texttt{infra/store} & Leases, optimistic concurrency, group-reservation saga, outbox, integrity audit \\
\texttt{infra/event\_bus} & Partitioned log with ordering, bounded retry, DLQ, lag \\
\texttt{infra/cache} & Versioned TTL cache, XFetch, single-flight \\
\texttt{sim/*} & Deterministic scenario generator and discrete-event harness \\
\texttt{tests/} & 36 tests covering exactness, monotonicity, conflict-freedom under thread contention, and degradation \\
\bottomrule
\end{tabular}
\end{center}

\subsection{Experimental design}

\begin{keybox}
342 incidents over 12 simulated hours against a fleet of 230 assets (112 ambulances, 40 fire units, 33 rescue teams, 4 helicopters, 33 hospitals) across 8 regions, with 50 injected environment disruptions and a 3.8-hour cyclone surge. Median of 3 seeds.

\medskip
Every policy is replayed against a \textbf{byte-identical} event stream
generated once from a fixed seed. Differences between rows are therefore
attributable to decision logic alone, not to the random number generator.
The ablation adds one mechanism at a time so that each contribution is
separately visible.
\end{keybox}

\subsection{Results}

\begin{center}\small\begin{tabular}{@{}lrrrrr@{}}
\toprule
Metric & Baseline & + Triage & + Re-opt & AEGIS & \textbf{Change} \\
\midrule
Mean first-response time (min) & 45.72 & 42.55 & 36.17 & 39.47 & \textcolor{good}{-13.7\%} \\
Median first-response (min) & 25.88 & 25.28 & 26.69 & 30.94 & \textcolor{bad}{+19.6\%} \\
p90 first-response (min) & 128.49 & 100.06 & 72.96 & 78.41 & \textcolor{good}{-39.0\%} \\
p99 first-response (min) & 208.19 & 205.01 & 146.69 & 155.81 & \textcolor{good}{-25.2\%} \\
Severity-weighted response (min) & 179.33 & 168.83 & 132.44 & 146.11 & \textcolor{good}{-18.5\%} \\
Reached within clinical window (\%) & 61.07 & 60.81 & 67.43 & 66.38 & \textcolor{good}{+8.7\%} \\
Never reached (\%) & 17.24 & 18.97 & 18.07 & 14.08 & \textcolor{good}{-18.3\%} \\
Fleet utilisation (\%) & 28.97 & 28.70 & 34.61 & 39.65 & \textcolor{good}{+36.9\%} \\
Assignment churn (\% of commits) & 0.00 & 0.00 & 2.96 & 6.10 & \textcolor{muted}{n/a} \\
\bottomrule
\end{tabular}\end{center}

\medskip
\begin{center}\small\begin{tabular}{@{}lrrrr@{}}
\toprule
Invariant / budget & Baseline & + Triage & + Re-opt & AEGIS \\
\midrule
Resource conflicts detected & 0 & 0 & 0 & 0 \\
Double-booking attempts blocked & 0 & 0 & 0 & 0 \\
Tier-1 admission p99 (ms) & 1.40 & 1.36 & 2.01 & 2.45 \\
Tier-2 epoch solve p99 (ms) & nan & nan & 38.70 & 38.53 \\
Travel-time cache hit rate (\%) & 77.4 & 78.3 & 98.2 & 97.1 \\
Degraded travel-time estimates served & 66 & 67 & 1,380 & 2,236 \\
Circuit-breaker trips & 12 & 14 & 30 & 30 \\
Events dead-lettered & 0 & 0 & 0 & 0 \\
Stranded leases auto-reclaimed & 0 & 0 & 0 & 0 \\
\bottomrule
\end{tabular}\end{center}

\subsection{Reading the results honestly}

\begin{itemize}
\item \textbf{The tail moves most.} Mean response improves 13.7\%; p90 improves
      39.0\% and p99 improves 25.2\%. That ordering is the expected signature
      of better allocation: the naive dispatcher handles easy cases adequately
      and fails badly on the contended ones, which is precisely where the exact
      solver earns its cost.

\item \textbf{The median gets worse, and that is the design working.} Median
      response rises from 25.9 to 30.9 minutes while every tail percentile
      falls sharply. This is redistribution, not regression: the baseline
      serves whoever is nearest, which flatters the median by handling easy
      calls first and abandoning contended ones. AEGIS deliberately makes a
      routine call wait so that a catastrophic one is reached at all
      \textemdash{} which is what \emph{triage} means. The severity-weighted
      mean, the metric that reflects that intent, improves 18.5\%. We report
      the median because hiding it would misrepresent the trade.

\item \textbf{Triage alone is not enough.} Adding the scoring model improves
      the tail (p90 falls 22\%) but slightly \emph{worsens} the in-window rate
      and the never-reached rate. Prioritising a queue does not create
      capacity; it only reorders who waits. The large gains arrive with
      re-optimisation, which is what actually reallocates.

\item \textbf{Coverage awareness trades speed for reach.} The full
      configuration is worse than the re-optimisation-only configuration on
      mean response (39.5 vs 36.2 minutes) and marginally on in-window rate
      (66.4\% vs 67.4\%), and clearly better on incidents never reached at all
      (14.1\% vs 18.1\%) and utilisation (39.7\% vs 34.6\%). Holding a reserve
      costs average speed and buys coverage for the calls that would otherwise
      get nothing. Whether that is the right trade is a policy question for
      clinicians and regulators, not an engineering one \textemdash{} which is
      exactly why the coverage weight $\theta$ is a configurable parameter
      with a measurable effect rather than a hard-coded rule. A service
      optimising a published mean-response target would set it lower; a
      service accountable for unreached calls would set it higher.

\item \textbf{Utilisation rises 37\%} while churn stays at 6.1\% of commits.
      The extra utilisation is real work delivered, not thrash: the engine
      committed 1{,}491 assignments against the baseline's 738 over the same
      stream, roughly double the unit-responses.

\item \textbf{Zero resource conflicts} in every configuration and every seed,
      corroborated by the multi-threaded contention test.

\item \textbf{Degraded operation was exercised.} The full configuration served
      2{,}236 travel-time estimates through the great-circle fallback during
      the injected routing outage without a single failed dispatch, while
      sustaining a 97\% cache hit rate. The baseline's much lower degraded
      count (66) simply reflects that it makes far fewer estimates, not that
      it is more resilient.
\end{itemize}

\subsection{Threats to validity}

Stated plainly, because a result whose limits are hidden is not a result.

\begin{itemize}
\item The simulator is our own model of the world; travel times are
      distance-based with detour and hazard multipliers rather than drawn from
      a real road network. Absolute minutes are therefore not calibrated to
      any real service, and only the \emph{relative} comparison is meaningful.
\item The baseline is a competent nearest-available dispatcher with
      priority-blind FIFO backfill. Real dispatch centres also have
      experienced humans; the comparison is against an automated policy, not
      against expert practice.
\item Three seeds establish that the direction of the effect is stable, not
      the precision of its magnitude.
\item Service-time and bed-occupancy distributions are plausible but not
      empirically fitted.
\end{itemize}

% ======================================================== limitations =====
\section{Limitations and roadmap}

\subsection{Known limitations of the current design}

\begin{itemize}
\item \textbf{No multi-stop sequencing.} A unit serves one incident, then
      returns to the pool. Real ambulances sometimes chain calls. This is a
      vehicle-routing extension, and the LNS layer is where it belongs.
\item \textbf{Cross-region borrowing is greedy.} Shards are solved
      sequentially with a claimed-resource set. A hierarchical coordinator
      solving a coarse national transfer problem would be better and is the
      natural next step.
\item \textbf{The demand forecast is deliberately simple} \textemdash{} an
      exponentially-weighted rate with a self-exciting term. A properly fitted
      spatio-temporal Hawkes process or a learned model would forecast better,
      and the interface is designed so it can be swapped without touching the
      optimiser.
\item \textbf{Coverage weight is uncalibrated.} Its effect is measured but the
      right value is a policy decision requiring clinical input.
\end{itemize}

\subsection{Roadmap}

\begin{enumerate}
\item Real road-network integration (OSRM with live traffic) and calibration
      against historical dispatch records.
\item Hierarchical cross-region coordinator with a national transfer problem.
\item Learned feature estimators upstream of the triage aggregator: predicted
      casualty count, predicted service duration, escalation probability.
\item Multi-stop sequencing and dynamic hospital-destination optimisation.
\item Shadow-mode policy evaluation pipeline as a standing capability, so
      every weight change ships with replayed evidence.
\end{enumerate}

% ========================================================== appendix ======
\appendix
\section{Appendix: an actual decision record}

Produced by the reference implementation during the evaluation run. This is the
artefact a dispatcher sees and an inquiry replays.

\begin{center}\footnotesize
\begin{tabular}{@{}lp{80mm}@{}}
\toprule
\multicolumn{2}{@{}l}{\textbf{Assignment} \texttt{ASG-0000042} \textemdash{} incident \texttt{INC-0000018}, unit \texttt{RES-0000117}} \\
\midrule
Tier & \texttt{T1\_admission} \\
Candidates considered & 36 \\
Candidates feasible & 9 \\
Rejected & \texttt{not\_available}: 21, \texttt{beyond\_reach\_radius}: 4, \texttt{capability\_below\_minimum}: 2 \\
\addlinespace
\multicolumn{2}{@{}l}{\emph{Why this incident ranked where it did}} \\
Priority & 71.4 \;(\texttt{triage-v1.3.0}) \\
\quad severity & 24.6 \quad (level 4) \\
\quad urgency & 16.1 \quad (28-minute window, 6 minutes elapsed) \\
\quad scale & 12.9 \quad (34 people affected) \\
\quad hazard & 9.0 \quad (structural collapse) \\
\quad vulnerability & 6.1 \\
\quad escalation & 2.7 \quad (aging, bounded at 12) \\
\addlinespace
\multicolumn{2}{@{}l}{\emph{Why this unit}} \\
ETA & 11.4 minutes \\
\quad travel term & $-63.4$ \\
\quad lateness term & $0.0$ \quad (arrives 16.6 minutes inside the window) \\
\quad over-qualification & $+3.0$ \quad (level-2 unit on a level-1 slot) \\
\quad scarcity & $0.0$ \quad (ambulance) \\
\quad coverage & $0.0$ \quad (region retains 6 free units) \\
\quad \textbf{total} & $\mathbf{-60.4}$ \\
Runner-up cost & $-48.1$ \quad (unit 4.1 minutes further away) \\
\bottomrule
\end{tabular}
\end{center}

\section{Appendix: reproducing these results}

\begin{center}\footnotesize
\begin{tabular}{@{}ll@{}}
\toprule
\textbf{Command} & \textbf{Effect} \\
\midrule
\texttt{python tests/test\_all.py} & Runs all 36 tests; no dependencies required \\
\texttt{python run\_simulation.py} & Full ablation, 12 simulated hours \\
\texttt{python run\_simulation.py -{}-seeds 3} & Repeats across seeds, reports medians \\
\texttt{python run\_simulation.py -{}-policy aegis} & Single policy, writes sample decision records \\
\texttt{python make\_tables.py} & Regenerates every table in this document from \texttt{artifacts/results.json} \\
\bottomrule
\end{tabular}
\end{center}

\noindent No number in this document was transcribed by hand; all tables are
generated from the results file produced by the simulation run.

\begin{keybox}
\textbf{A note on determinism.} In production a re-optimisation epoch is
bounded by wall-clock time, because that is what a latency budget means. That
makes the number of shards refreshed per epoch depend on machine load, and
therefore makes results irreproducible. Evaluation runs instead bound the epoch
by a \emph{deterministic work budget} (shards per epoch), configured by
\texttt{OptimizerConfig.deterministic\_budget}. This was found by running the
same evaluation twice and getting different answers; it is the difference
between a measurement and an anecdote.
\end{keybox}

\vspace{6mm}
\begin{keybox}
\textbf{Closing note on method.} Four substantive bugs were found by building
the system rather than by describing it: a feasibility cutoff missing from an
exact matcher that consequently dispatched units from hundreds of kilometres
away; a coverage penalty charged against already-dispatched units, which made
every busy vehicle look expensive relative to an idle one; slot-index rather
than demand-level comparison, which reported interchangeable columns as
displacements; and stale position estimates for units already in motion. Each
one individually made the "optimised" system perform \emph{worse} than the
naive baseline while every architecture diagram still looked correct. That is
the argument for a running reference implementation.
\end{keybox}

\end{document}
